\documentclass[12pt,a4paper,oneside]{book}

% border setting
\usepackage[ top=2.5cm,bottom=2.5cm,left=2.5cm,right=2.5cm ]{ geometry }

% The default for LaTeX is to have no indent after sectional headings, like \chapter and \section. ()
\usepackage{indentfirst}

% http://tex.stackexchange.com/questions/28333/continuous-v-per-chapter-section-numbering-of-figures-tables-and-other-docume
\usepackage{chngcntr}
\counterwithout{figure}{chapter}
\counterwithout{table}{chapter}

% This prevents placing floats before a section
\usepackage[section]{placeins}
\let\Oldsubsection\subsection
\renewcommand{\subsection}{\FloatBarrier\Oldsubsection}

% source code hightlighting
\usepackage{listings}
\lstset{
  numbers=left,
  stepnumber=1,
  firstnumber=1,
  captionpos=b,
  tabsize=2,
  basicstyle=\small,
  numberfirstline=true
}

% setting the page number to footer
\usepackage{fancyhdr}
\fancyhf{}
\cfoot{\thepage}
\pagestyle{fancy}
% no header and footer bar
\renewcommand{\headrulewidth}{0pt}
\renewcommand{\footrulewidth}{0pt}

% setup bibliography
\usepackage[sorting=none,backend=bibtex]{biblatex}
\addbibresource{reference.bib}

% line height setting
\linespread{1.5}
\usepackage{setspace}

% Graphics settings
\usepackage{graphicx}
\graphicspath{ {./figures/} }

\usepackage{background}
\newcommand\DeactivateBG{\backgroundsetup{contents={}}}
\newcommand\ActivateBG{
  \backgroundsetup{
      contents={\includegraphics[]{logo.jpg}},
      scale=1,
      opacity=0.4,
      angle=0
  }
}



% do not place figure at the middle of a empty page
\makeatletter
\setlength{\@fptop}{0pt}
\makeatother

% Chinese typesetting
\usepackage{xeCJK}
%\setCJKmainfont{SourceHanSansTW-Light.otf}
\setCJKmainfont[
  BoldFont={SourceHanSansTW-Normal.otf},
  ItalicFont={SourceHanSansTW-Light.otf}
]{SourceHanSansTW-Light.otf}
\newcommand{\myHuge}[1]{\fontsize{40}{50} #1}

\newcommand{\chineseTitle}{中文題目}
\newcommand{\englishTitle}{英文題目}

\newcommand{\studentCnName}{葉家郡}
\newcommand{\studentEnName}{Jia-Jun Yeh}
\newcommand{\advisorCnName}{吳小松}
\newcommand{\advisorEnName}{Xiao-Sung Wu}

\usepackage{pdfpages}

\begin{document}

\pagenumbering{gobble} % disabling page numbering

\ActivateBG

\begin{titlepage}
  \begin{center}
    \myHuge \textbf{國立交通大學}    \\[0.25cm]

    \Huge \textbf{資訊科學與工程研究所} \\[0.25cm]

    \Huge \textbf{碩士論文} \\[1cm]
	
    \LARGE \chineseTitle{} \\[0.5cm]

    \LARGE \englishTitle{} \\
  \end{center}

  \vspace{\fill}

  \begin{tabular}{c l}
    {\makebox[8em][s]{\LARGE 研究生}}   & \LARGE ：\studentCnName{}    \\[0.5cm]
    {\makebox[8em][s]{\LARGE 指導教授}} & \LARGE ：\advisorCnName{} \hspace{0.1cm} 教授\\
  \end{tabular}

  \vspace{3cm}

  \begin{center}
    {\LARGE 中華民國  105 年 7 月}
  \end{center}
\end{titlepage}

\begin{titlepage}
  \begin{center}
    \LARGE \chineseTitle{}  \\
    \LARGE \englishTitle{}  \\[1.5cm]
  
    \Large
    \begin{tabular}{c l c c l}
    研究生   & ：\studentCnName{} & \hspace{3cm}  & Student  & ：\studentEnName{} \\
    指導教授 & ：\advisorCnName{} & \hspace{3cm}  & Advisor  & ：\advisorEnName{}\\
    \end{tabular}
    \\[1.5cm]
    國立交通大學 \\
    資訊科學與工程研究所 \\
    碩士論文 \\[1cm]
	
    \begin{singlespace}
    A Thesis Submitted to Institute of Computer Science and Engineering College of Computer Science National Chiao Tung University in Partial Fulfillment of the Requirements for the Degree of Master in Computer and Information Science \\
    July 2016 \\
    \studentEnName{}, Taiwan \\
    \end{singlespace}

  \end{center}

  \vspace{\fill}

  \begin{center}
    {\LARGE 中華民國  105 年 7 月}
  \end{center}
\end{titlepage}

\DeactivateBG

% 放審定書和授權書
% \includepdf[pages={1-2}]{auth.pdf}

\ActivateBG

\begin{titlepage}
  \begin{center}
  	\LARGE 
    \begin{singlespace}    
      \textbf{\chineseTitle{}} \\[0.5cm]
    \end{singlespace}
    
    \begin{singlespace}    
    \begin{tabular}{r l}
    	學生     & ：\studentCnName{}  \\
        指導教授  & ：\advisorCnName{} \hspace{0.1cm} 教授 \\[0.5cm]
    \end{tabular}
    \end{singlespace}

    國立交通大學資訊科學與工程研究所碩士班 \\[0.5cm]
    \makebox[4em][s]{摘要} \\[0.5cm]
    	
  \end{center}
  \normalsize 
  \hspace{0.6cm} 在程式自動化測試與分析上，符號執行(symbolic execution)是目前經常被使用的一種方法，由於符號執行會紀錄並模擬出程式執行時的所有可能路徑，其數量會以指數的數量級成長，最終耗盡所有運算資源，這個問題被稱為路徑爆炸問題(path explosion problem)；因此我們需要在有限的資源內採取某些策略來優先計算較有價值的路徑，在本篇論文中我們提出使用以蒙地卡羅搜尋樹為基礎的搜尋策略來解決這個問題，並比較它與其他傳統策略如深度優先搜尋(DFS)、廣度優先搜尋(BFS)的效率。
  \\[0.7cm]
  關鍵字：Monte Carlo Tree Search(MCTS), Upper Confidence Bounds for Trees (UCT), symbolic execution
\end{titlepage}
\begin{titlepage}
  \begin{center}
    \LARGE
    \begin{singlespace}
  	 \textbf{\englishTitle{}} \\[0.5cm]
    \end{singlespace}
    
    \begin{singlespace}
    \begin{tabular}{r l}
    	Student     & : \studentEnName{}  \\
        Advisor  & : Dr. \advisorEnName{} \\[0.5cm]
    \end{tabular}
    \end{singlespace}
	
    \begin{singlespace}
    Institute of Computer Science and Engineering National Chiao Tung University\\[0.5cm]
    \end{singlespace}
    \textbf{ABSTRACT} \\[0.5cm]
    	
  \end{center}
  \normalsize 
  \hspace{0.6cm} Symbolic execution is a technology that is often used today in program automation testing and analysis. Since symbol execution traces and simulates all possible paths when a program executes, its number grows exponentially. This problem is called the path explosion problem. Therefore, we need to take some strategies within the limited resources to give priority to the more valuable path. In this paper, we propose to use The Carlo search tree-based search strategy solves this problem and compares it with other classic strategies such as depth-first search (DFS) and breadth-first search (BFS).
  \\[0.7cm]
  Keywords：Monte Carlo Tree Search(MCTS), Upper Confidence Bounds for Trees (UCT), symbolic execution
\end{titlepage}

\tableofcontents
\listoffigures
\listoftables

% 9/30
\chapter{Introduction} \pagenumbering{arabic} % enabling page numbering

在現今的軟體開發中，程式碼數量動輒數萬行，系統的複雜度也越來越高，傳統驗證程式正確性的方式如unit testing、code review等等...受限於人工而相當有限；在硬體計算能力突飛猛進的現代，自動化測試的方式又逐漸成為顯學，如fuzzing、symbolic execution...能夠自動尋找程式中可能的漏洞，其中symbolic execution是一種模擬執行的方法，它將程式的使用者輸入視為符號，並把程式執行過程和分支條件轉換為限制式，藉由求解限制式來獲得欲執行該路徑所需的使用者輸入為何；由於symbolic execution在遇到分支時會複製出一條新的路徑，兩條路徑分別探索執行if時和執行else時的狀態，因此路徑數量會以指數的數量級成長，造成探索路徑需要花費巨大的運算資源，這個問題被稱為path explosion problem，這篇論文欲以蒙地卡羅樹搜尋演算法來找出探索價值較高的路徑，使得symbolic execution能使用較短的時間內獲得更高的程式執行覆蓋率。

蒙地卡羅樹搜尋被廣泛的運用在遊戲人工智慧中，例如西洋棋、黑白棋、圍棋等等的棋盤遊戲，在2016年AlphaGo(一個結合蒙地卡羅樹搜尋和深度學習的圍棋AI程式)擊敗世界棋王後，更是一度引起相當多的討論；此演算法主要的精神在於對一棵樹，選擇子結點並嘗試用模擬的方式估計該結點的價值。和symbolic execution相當類似的地方在於同為需要探索一棵樹，而且要避免探索一些我們不感興趣或是沒有必要探索的路徑，例如無窮迴圈的狀況；我們認為結合蒙地卡羅樹搜尋做為symbolic execution挑選路徑執行的演算法，相較於傳統的深度優先搜尋法或廣度優先搜尋法，能有較佳的搜尋效率。

% \cite{thesis}.

% 9/30
\chapter{Background}

\section{Symbolic execution}

\section{Path Exploration}

\section{Monte Carlo tree search}

\chapter{Design}
  
\chapter{Implementation}

\chapter{Evaluation}

% 9/30
\chapter{Related Work}

\chapter{Conclusion}

\newpage

\printbibliography[title={References}]

\end{document}
